\documentclass[12pt]{article}
\usepackage{amsmath,amssymb,amsthm}
\usepackage{mathtools}
\usepackage{geometry}
\geometry{margin=1in}

% Theorem environments
\newtheorem{definition}{Definition}[section]
\newtheorem{lemma}[definition]{Lemma}
\newtheorem{theorem}[definition]{Theorem}
\newtheorem{corollary}[definition]{Corollary}
\newtheorem{proposition}[definition]{Proposition}
\newtheorem{remark}[definition]{Remark}

% Custom commands
\newcommand{\B}{\mathbb{B}}
\newcommand{\C}{\mathbb{C}}
\newcommand{\R}{\mathbb{R}}
\renewcommand{\H}{\mathbb{H}}
\newcommand{\Tr}{\mathrm{Tr}}
\newcommand{\re}{\mathrm{Re}}
\newcommand{\im}{\mathrm{Im}}

\title{Step 3: Einstein Field Equations with Matter from UBT\\
\large $G_{\mu\nu} = 8\pi G\,T_{\mu\nu}$ Derived, Not Asserted}
\author{UBT Theory Development}
\date{2026}

\begin{document}

\maketitle

\begin{abstract}
We derive the full Einstein equations with matter $G_{\mu\nu} = 8\pi G\,T_{\mu\nu}$
from the combined UBT action $S_{\mathrm{total}}[g,\Theta]
= S_{\mathrm{EH}}[g] + S_{\mathrm{matter}}[\Theta, g]$.
Varying $S_{\mathrm{total}}$ with respect to $g^{\mu\nu}$ and using the
results of Steps~1--2, we show that the stationary condition reproduces
Einstein's equations with $T_{\mu\nu} = \mathrm{Re}(\mathcal{T}_{\mu\nu})$.
We also address energy-momentum conservation $\nabla^\mu T_{\mu\nu} = 0$
via the Bianchi identity and Noether's theorem.
This closes item~II.9 of \texttt{THEORY\_INCONSISTENCIES\_REPORT.md}.
\end{abstract}

\tableofcontents

% ------------------------------------------------------------
\section{The Total UBT Action}

\subsection{Einstein--Hilbert Term}

\begin{definition}[Einstein--Hilbert action]
\label{def:EH}
\begin{equation}
  \label{eq:S_EH}
  S_{\mathrm{EH}}[g] = \frac{1}{16\pi G}\int d^4x\,\sqrt{-g}\; R,
\end{equation}
where $R = g^{\mu\nu}R_{\mu\nu}$ is the Ricci scalar and
$G$ is Newton's gravitational constant.
\end{definition}

\subsection{Matter Action}

From \texttt{appendix\_AA\_theta\_action.tex} and Steps~1--2:
\begin{equation}
  \label{eq:S_matter}
  S_{\mathrm{matter}}[\Theta, g]
  = S_{\mathrm{kin}} + S_{\mathrm{pot}} + S_{\mathrm{gauge}}.
\end{equation}

\subsection{Combined Action}

\begin{definition}[Full UBT action]
\label{def:S_total}
\begin{equation}
  \label{eq:S_total}
  S_{\mathrm{total}}[g, \Theta]
  = S_{\mathrm{EH}}[g] + S_{\mathrm{matter}}[\Theta, g].
\end{equation}
The dynamical variables are the metric $g^{\mu\nu}$ and the biquaternionic
field $\Theta$.
\end{definition}

% ------------------------------------------------------------
\section{Variation with Respect to the Metric}

\subsection{Palatini Identity for the Einstein--Hilbert Term}

\begin{lemma}[Hilbert variation of $S_{\mathrm{EH}}$]
\label{lem:EH_var}
\begin{equation}
  \label{eq:var_EH}
  \frac{\delta S_{\mathrm{EH}}}{\delta g^{\mu\nu}}
  = \frac{\sqrt{-g}}{16\pi G}
    \Bigl(R_{\mu\nu} - \tfrac{1}{2}g_{\mu\nu}R\Bigr)
  = \frac{\sqrt{-g}}{16\pi G}\,G_{\mu\nu},
\end{equation}
where $G_{\mu\nu} = R_{\mu\nu} - \tfrac{1}{2}g_{\mu\nu}R$ is the Einstein tensor.
\end{lemma}

\begin{proof}
  Standard derivation.  Using $\delta(\sqrt{-g}R) = \sqrt{-g}
  (R_{\mu\nu} - \tfrac{1}{2}g_{\mu\nu}R)\delta g^{\mu\nu}
  + \sqrt{-g}\,\nabla_\mu V^\mu$ where $V^\mu$ is the Palatini boundary
  vector, and discarding the total derivative term under the integral sign.
  See Wald~\cite{wald}, Appendix~E.
\end{proof}

\subsection{Variation of the Matter Action}

From Definition~2.1 of \texttt{step1\_metric\_variation.tex}:
\begin{equation}
  \label{eq:var_matter}
  \frac{\delta S_{\mathrm{matter}}}{\delta g^{\mu\nu}}
  = -\frac{\sqrt{-g}}{2}\,T_{\mu\nu},
\end{equation}
where $T_{\mu\nu} = \re(\mathcal{T}_{\mu\nu})$ is the stress-energy tensor
derived in Steps~1--2.

\subsection{Combined Stationarity Condition}

Setting $\delta S_{\mathrm{total}}/\delta g^{\mu\nu} = 0$:
\begin{equation}
  \frac{\delta S_{\mathrm{EH}}}{\delta g^{\mu\nu}}
  + \frac{\delta S_{\mathrm{matter}}}{\delta g^{\mu\nu}} = 0.
\end{equation}
Substituting~\eqref{eq:var_EH} and~\eqref{eq:var_matter}:
\begin{equation}
  \frac{\sqrt{-g}}{16\pi G}\,G_{\mu\nu}
  - \frac{\sqrt{-g}}{2}\,T_{\mu\nu} = 0.
\end{equation}
Dividing by $\sqrt{-g}/2$:
\begin{equation}
  \frac{1}{8\pi G}\,G_{\mu\nu} = T_{\mu\nu}.
\end{equation}

% ------------------------------------------------------------
\section{Einstein Field Equations with Matter}

\begin{theorem}[Einstein equations from UBT]
\label{thm:einstein}
Stationarity of the total UBT action~\eqref{eq:S_total} with respect to
the metric $g^{\mu\nu}$ yields:
\begin{equation}
  \label{eq:einstein_matter}
  \boxed{G_{\mu\nu} = 8\pi G\,T_{\mu\nu},}
\end{equation}
i.e.\ the full Einstein field equations \emph{with matter}, where
\begin{equation}
  T_{\mu\nu}
  = \re(\mathcal{T}_{\mu\nu})
  = \Tr\!\bigl[(\nabla_\mu\Theta)^\dagger\nabla_\nu\Theta\bigr]
    -\tfrac{1}{2}g_{\mu\nu}\Tr\!\bigl[(\nabla_\alpha\Theta)^\dagger
    \nabla^\alpha\Theta\bigr]
    - g_{\mu\nu}V(\Theta)
    + K\Bigl[\Tr[F_{\mu\alpha}F_\nu{}^\alpha]
      -\tfrac{1}{4}g_{\mu\nu}\Tr[F_{\alpha\beta}F^{\alpha\beta}]\Bigr].
\end{equation}
\end{theorem}

\begin{remark}[Vacuum limit]
  Setting $\Theta = 0$ and $A_\mu = 0$ (vacuum), all matter terms vanish
  and~\eqref{eq:einstein_matter} reduces to the vacuum Einstein equations
  $G_{\mu\nu} = 0$, as derived in \texttt{appendix\_R\_GR\_equivalence.tex}.
  $\checkmark$
\end{remark}

\begin{remark}[No new parameters]
  Equation~\eqref{eq:einstein_matter} introduces no new free parameters.
  The coupling constant $8\pi G$ is already present in $S_{\mathrm{EH}}$
  via Definition~\ref{def:EH}.
\end{remark}

% ------------------------------------------------------------
\section{Energy-Momentum Conservation}
\label{sec:conservation}

\subsection{Bianchi Identity Route}

\begin{theorem}[Automatic conservation from contracted Bianchi identity]
\label{thm:bianchi}
The contracted Bianchi identity $\nabla^\mu G_{\mu\nu} = 0$ holds
identically for any metric $g_{\mu\nu}$.  Combined with the Einstein
equations~\eqref{eq:einstein_matter}, it implies:
\begin{equation}
  \label{eq:conservation}
  \nabla^\mu T_{\mu\nu} = 0.
\end{equation}
This holds for \emph{any} $T_{\mu\nu}$ satisfying~\eqref{eq:einstein_matter},
regardless of the specific form of the matter action.
\end{theorem}

\begin{proof}
  From~\eqref{eq:einstein_matter}: $\nabla^\mu(8\pi G\,T_{\mu\nu})
  = \nabla^\mu G_{\mu\nu} = 0$ (contracted Bianchi).
  Since $8\pi G \neq 0$, we obtain $\nabla^\mu T_{\mu\nu} = 0$.
\end{proof}

\begin{remark}[No extra condition needed]
  Conservation~\eqref{eq:conservation} is \emph{automatic}; it does not
  require any additional assumption on $\Theta$.  It is a consequence of
  diffeomorphism invariance of the total action (see also the Noether
  route below).
\end{remark}

\subsection{Noether Route (Diffeomorphism Invariance)}

\begin{proposition}[Noether conservation]
\label{prop:noether}
The total action $S_{\mathrm{total}}$ is invariant under diffeomorphisms
$x^\mu \to x^\mu + \xi^\mu(x)$.  By Noether's second theorem applied to
the matter action:
\begin{equation}
  \nabla^\mu T_{\mu\nu} = 0
\end{equation}
holds as a \emph{Noether identity}, valid on the constraint surface
(when $\Theta$ satisfies its own Euler--Lagrange equations).
\end{proposition}

\begin{proof}[Proof sketch]
  Under $x^\mu \to x^\mu + \xi^\mu$, the metric transforms as
  $\delta g_{\mu\nu} = -\mathcal{L}_\xi g_{\mu\nu}
  = -\nabla_\mu\xi_\nu - \nabla_\nu\xi_\mu$.
  Invariance of $S_{\mathrm{matter}}$ gives:
  \begin{equation}
    0 = \delta S_{\mathrm{matter}}
      = -\int d^4x\,\sqrt{-g}\,T^{\mu\nu}\nabla_\mu\xi_\nu
      = \int d^4x\,\sqrt{-g}\,(\nabla^\mu T_{\mu\nu})\xi^\nu,
  \end{equation}
  after integration by parts (boundary terms vanish for $\xi^\mu$ with
  compact support).  Since $\xi^\nu$ is arbitrary,
  $\nabla^\mu T_{\mu\nu} = 0$.
\end{proof}

\begin{remark}[Conditions for conservation]
  Both routes give $\nabla^\mu T_{\mu\nu} = 0$:
  \begin{itemize}
    \item \textbf{Bianchi route:} requires only that $\Theta$ and $g$ jointly
          satisfy the full system of equations (Einstein + matter);
          no additional assumption.
    \item \textbf{Noether route:} requires that $\Theta$ satisfies the
          Euler--Lagrange equations (the $\Theta$ equation of motion from
          Appendix~AA, Theorem~4.1) \emph{and} that the action is
          diffeomorphism-invariant (which it is by construction).
  \end{itemize}
  In both cases no additional free parameters or conditions are needed
  beyond the existing UBT action principle.
\end{remark}

% ------------------------------------------------------------
\section{Summary and Connection to Existing Proofs}

\begin{corollary}[Closure of gap in \texttt{appendix\_R\_GR\_equivalence.tex}]
  \label{cor:gap_closed}
  The identification $T_{\mu\nu} = \re(\mathcal{T}_{\mu\nu})$ in
  \texttt{appendix\_R\_GR\_equivalence.tex} is now fully justified by the
  derivation chain:
  \begin{equation}
    S[\Theta]
    \xrightarrow{\delta/\delta g^{\mu\nu}\,(S_{\mathrm{matter}})}
    T_{\mu\nu}
    \xrightarrow{\delta/\delta g^{\mu\nu}\,(S_{\mathrm{EH}} + S_{\mathrm{matter}})}
    G_{\mu\nu} = 8\pi G\,T_{\mu\nu}.
  \end{equation}
\end{corollary}

Key results established in this chain:
\begin{enumerate}
  \item $T_{\mu\nu}$ is derived from $S[\Theta]$ by the Hilbert prescription
        (not asserted).
  \item $G_{\mu\nu} = 8\pi G\,T_{\mu\nu}$ follows from
        $\delta S_{\mathrm{total}}/\delta g^{\mu\nu} = 0$.
  \item $\nabla^\mu T_{\mu\nu} = 0$ is automatic (Bianchi + Noether).
  \item In the vacuum limit ($\Theta=0$), Einstein vacuum equations are
        recovered.
  \item In the QED limit, $T_{\mu\nu} \to T^{(\mathrm{EM})}_{\mu\nu}
        + T^{(\mathrm{Dirac})}_{\mu\nu}$ (Step~2, Proposition~5.1).
\end{enumerate}

The properties of $T_{\mu\nu}$ (realness, symmetry, divergence-free) are
verified in \texttt{step4\_T\_properties.tex}.

\begin{thebibliography}{99}
\bibitem{wald}    Wald, R.~M. (1984). \textit{General Relativity}.
                  University of Chicago Press.
\bibitem{misner}  Misner, C.~W., Thorne, K.~S., \& Wheeler, J.~A. (1973).
                  \textit{Gravitation}. Freeman.
\bibitem{step1}   UBT Theory Development.
                  \texttt{step1\_metric\_variation.tex}, 2026.
\bibitem{step2}   UBT Theory Development.
                  \texttt{step2\_biquat\_stress\_energy.tex}, 2026.
\bibitem{appendix_AA}
                  UBT Theory Development.
                  \texttt{appendix\_AA\_theta\_action.tex}, 2025.
\bibitem{appendix_R}
                  UBT Theory Development.
                  \texttt{appendix\_R\_GR\_equivalence.tex}, 2025.
\end{thebibliography}

\end{document}
